\chapter{Method Logic Trace}

This appendix records the reasoning chain behind the analysis choices.  It is
included to make the report auditable: a reader should be able to see not only
what number was quoted, but why that operation was used, what alternative it
guards against, and what interpretation is allowed.

\section{Global rules}

\begin{enumerate}
  \item Start from raw waveforms whenever a population count or selection claim
  matters.  Sorted or derived branches are useful diagnostics, but they can
  encode historical choices that are invisible to a later reader.
  \item Separate deterministic definitions from learned approximations.  If the
  label is exactly \(A>1000\) ADC, a model trained to predict that label is not
  discovering physics; it is approximating the cut.
  \item Compare methods on run-held-out data.  Run period, beam current,
  topology, and calibration are correlated, so random row splits can overstate
  transfer performance.
  \item Accept ML only against a strong traditional comparator on the same
  target, population, split, and metric.  A neural model beating a weak baseline
  does not justify replacing a transparent physics model.
  \item Distinguish construction truth, electronics closure, simulation truth,
  and real-event physics truth.  Injected overlaps, duplicate readout, and
  GEANT4 each answer different questions.
\end{enumerate}

\section{Step-by-step chain}

\begin{longtable}{@{}p{0.06\linewidth}p{0.20\linewidth}p{0.24\linewidth}p{0.24\linewidth}p{0.18\linewidth}@{}}
\toprule
Step & Input & Operation & Reason & Interpretation \\
\midrule
\endhead
1 & Raw B-stack \texttt{HRDv} & Read the 18-sample waveforms for physical even B channels & This is the least ambiguous source for a pulse-count claim & Counts are raw-data anchored \\
2 & Samples 0--3 & Compute \(b=\mathrm{median}(w_0,w_1,w_2,w_3)\) & The median is robust to one bad pre-pulse sample & This is a selection baseline, not final pedestal truth \\
3 & Baseline-subtracted waveform & Compute \(A=\max_j(w_j-b)\) & Peak amplitude is the historical selection scalar & The selected population is deterministic \\
4 & \(A\), stave id & Keep B2/B4/B6/B8 pulses with \(A>1000\) ADC & This defines the reproducible B-stack population & All downstream claims inherit this gate \\
5 & Selected pulses & Build stave- and amplitude-binned templates & Pulse shape changes with stave and amplitude & Template quality is a covariate \\
6 & Leading edge & Compute CFD20 seed time & A fractional crossing gives an interpretable local seed & The seed is a software primitive \\
7 & Template and derivative & Fit sub-sample template phase & Timing needs a local shape model and sample masking & Raw time is model-estimated \\
8 & Amplitude and run structure & Apply timewalk and offsets & Leading-edge time depends on amplitude and calibration period & Analytic timewalk explains most timing gain \\
9 & Corrected stave times & Form inter-stave residuals & Same-particle residuals avoid needing external absolute time truth & B4/B6/B8 define the clean timing reference \\
10 & Residual distribution & Quote \(\sigma_{68}\), core fits, and tails when available & A single Gaussian core hides topology tails & The metric must be named with every claim \\
11 & B2 residuals & Compare topology covariance & B2 has terminal and overlap-dominated populations & B2 broadening is topology-local \\
12 & Occupancy target & Compute \(R_{\max}=\mu_{\max}/\tau_{\mathrm{eff}}\) & Rate limits depend directly on live-time & The 4.22 MHz value assumes 90 ns \\
13 & Waveform tails & Measure live-time from data & The assumed 90 ns window must be tested & The 10-percent estimate implies about 3.05 MHz \\
14 & Current pairs & Compare high-current excess after baseline structure & Real beam pile-up should increase with current & Quote excess, not raw score \\
15 & Injected overlaps & Compare bounded template and ML recovery & Synthetic truth allows controlled method ranking & ML RMS wins require failure-rate gates \\
16 & Waveform matrices & Compare PCA and autoencoder latents & Tests whether pulse shape is low-dimensional & Nonlinearity helps mainly in compact latents \\
17 & Duplicate readout & Predict amplitude and charge closure & The target is independent electronics information & This is not absolute energy truth \\
18 & Artificial clips & Recover unclipped amplitude & Construction truth is known & ML wins artificial saturation closure \\
19 & Pretrigger samples & Validate pedestal estimators & A positivity constraint can pass by construction & Adaptive lowering is a pathology diagnostic \\
20 & GEANT4 truth & Benchmark energy models & Real HRD data lack event-level energy and PID truth & GEANT4 anchors energy interpretation \\
\bottomrule
\end{longtable}

\section{How the conclusions follow}

The selection conclusion follows because the gate contains no trained parameter:
for fixed input bytes and fixed implementation, the row count is a reproducible
integer.  The uncertainty is therefore provenance and implementation risk, not
binomial counting error.

The timing conclusion follows because a stronger analytic amplitude-timewalk
model closes most of the initial ML advantage.  ML residual correction remains a
useful diagnostic, but it is not adopted as the default timing model unless it
wins against the analytic comparator on the same run-held-out benchmark with a
larger and calibrated margin.

The pile-up-rate conclusion follows from the occupancy equation.  The previous
headline was conditional on \(\tau_{\mathrm{eff}}=90\) ns.  Once the waveform
tail measurement gives a longer live-time, the same occupancy requirement
implies a lower admissible rate.  ML classifiers diagnose waveform structure,
but they do not replace the live-time term in the rate equation.

The ML conclusion is task-specific.  ML wins where the target is independent and
shape contains recoverable information, such as duplicate-readout closure and
artificial saturation.  It is rejected or kept diagnostic where the label is a
restatement of the inputs, where leakage risk is high, or where a transparent
physics model already captures the dominant structure.

The energy and PID conclusion follows from missing truth.  HRD waveform data can
support sample-level enrichment and closure tests, but they do not contain
event-level particle identity or deposited-energy truth.  GEANT4 is therefore
the required bridge for physics-facing energy and PID claims.

\section{Future edit checklist}

Any future change to a headline number must update the source study reference,
the Markdown summary, the LaTeX chapter, the visual figure or table, the
validation command, and the caveat that bounds the interpretation.  This prevents
a table from changing while the reasoning chain remains stale.
